\chapter{Combustible y control ambiental}
\label{cap:mm-combustible}

\opening{Un sistema por gravedad, sin bombas: dos tanques, una válvula, un gascolador y unos respiraderos que no deben obstruirse.}

\section{Sistema de combustible: descripción}
\fuente{MM, capítulo 12}

El mantenimiento depende del tipo de tanque: las aeronaves de producción actual tienen \textbf{vejigas flexibles dentro de envolventes de aluminio}; las anteriores, \textbf{tanques totalmente de aluminio}. Identifique el tipo antes de trabajar. El carburador y los accesorios de combustible del motor se mantienen según los datos de su fabricante.

El sistema es \textbf{por gravedad} (sin bombas) e incluye tanque principal, tanque auxiliar (opcional en los de aluminio), válvula de corte detrás del asiento izquierdo y filtro (gascolador). Los respiraderos de los tanques están dentro del carenado del mástil.

Hay válvulas de drenaje tipo émbolo en el gascolador (lado inferior izquierdo del cortafuego) y en el sumidero de cada tanque; las del tanque auxiliar y del principal de vejiga están dentro de la puerta del capó, bajo el tanque auxiliar, con tubos plásticos que descargan por la borda. En los tanques principales de aluminio, el drenaje está en el lado izquierdo del fuselaje y se abre empujando el tubo plástico. En helicópteros nuevos hay un tubo de vidrio guardado en la puerta superior trasera del capó para recoger muestras. \textbf{Se debe tomar muestra en los tres puntos antes del primer vuelo del día y después de cada recarga}, para verificar ausencia de contaminación y el grado correcto.

Los indicadores son eléctricos, con transmisores de flotador; con ``E'' los tanques están vacíos salvo una pequeña cantidad no utilizable. La luz de bajo combustible la acciona un emisor separado en el fondo del tanque principal. El tanque auxiliar está interconectado con el principal y algo más alto, por lo que se vacía primero. La válvula de corte controla el flujo de ambos.

\begin{advertencia}
Los vapores de combustible son explosivos: no use herramientas eléctricas cerca de un sistema abierto. Las mangueras flexibles se doblan fácilmente; tape temporalmente todo racor abierto.
\end{advertencia}

\section{Tanques}

\subsection{Tanque principal}

Para retirarlo se desabastece la aeronave, se retiran carenado del mástil, carenado del cono (desconectando antena y masa del ELT si lo hay), canal B423-1 y revestimientos B429, respaldo izquierdo; se desconectan respiraderos, correa A713-1, cableado del emisor de cantidad y del interruptor de bajo combustible, y las mangueras (en vejiga, quedan unidas al tanque) o líneas (aluminio).

Puntos clave de la instalación:
\begin{itemize}
\item La tira vertical del cortafuego horizontal debe quedar entre el ángulo D249-4 del tanque y su revestimiento exterior.
\item Verifique eléctricamente el interruptor de bajo combustible y el emisor de cantidad antes de instalar.
\item Tuercas de mangueras y líneas con torque especial y marca. Mangueras D205-33 y D205-34 con bridas plásticas y \textbf{0,25~in} mínimo de holgura con bastidor, cortafuego y soporte.
\item Respiraderos con abrazaderas apretadas o dos vueltas de alambre de 0,032~in.
\item Verifique la indicación de cantidad y la ausencia de fugas; las franjas de alineación de la tapa (SB-107) y las calcomanías D682-44; si se tocó, haga prueba de fugas pitot-estática.
\end{itemize}

\subsection{Tanque auxiliar}

Procedimiento similar, abriendo además el lado derecho del carenado del mástil y la puerta B430, retirando el respaldo derecho, la abrazadera de la línea pitot-estática y el cable del freno rotor de la palanca B112-1 (cadena unida con termocontraíble B158-6). La manguera D205-33 debe mantener 0,25~in de holgura con bastidor y cortafuego.

\subsection{Cambio de vejiga}

\begin{precaucion}
La vejiga (B028-1 principal, B028-2 auxiliar) debe estar \textbf{por encima de 65~°F} antes de retirarla, instalarla o flexionarla. Durante el desarme, encinte progresivamente las aberturas para evitar objetos extraños.
\end{precaucion}

Se retiran emisor de cantidad, tapa inferior (con el interruptor de bajo combustible en el principal; con filtro, sifón y unión en el auxiliar), tornillos de respiraderos y boca de llenado, filtro y uniones (principal), mamparo, canal y rigidizador, y se despega el velcro. Al armar: tornillos con antiagarrotante A257-9; juntas tóricas nuevas lubricadas con grasa A257-6; torque especial en filtro y uniones. \textbf{Prueba de presión}: tapando respiraderos y racores, presurice la vejiga a \textbf{no más de 1~psi} y busque fugas con agua jabonosa en todas las costuras.

\subsection{Tapas de combustible}

\begin{advertencia}
Las franjas de alineación de las tapas son obligatorias por reglamento y permiten al piloto verificar que la tapa está bien cerrada; una tapa mal asegurada puede salirse en vuelo.
\end{advertencia}

La tapa A689 lleva dos franjas a 180° y el cuello A250 una franja hacia fuera (anchos de 0,20--0,30~in y 0,30--0,40~in). Si se deterioran: limpiar con acetona y pintar con Imron negro D257-7 en aeronaves claras o blanco D257-5 en oscuras. La calcomanía D682-44 se alinea con las franjas.

\section{Respiraderos}

Las válvulas antivuelco D251-1 tienen una bola de latón: normalmente el aire pasa alrededor de ella desde los espacios de expansión; si hay sobrellenado o chapoteo, la bola flota y sella el respiradero, y con la aeronave invertida lo sella para minimizar derrames. No tienen empaquetaduras internas. Se instalan respetando las marcas de orientación.

\textbf{Verificación del sistema de respiraderos}: con las tapas puestas, conecte una manguera temporal a uno de los tubos de respiradero A731-4 y sople (\textbf{nunca con aire comprimido}): el aire debe salir por el otro. Si no, elimine la obstrucción.

\section{Válvulas de drenaje}

Válvula A761-1: junta tórica nueva en el asiento del vástago, rosca con sellador B270-6 liviano, torque estándar y marca; tubo A729 asegurado con dos vueltas de alambre de 0,032~in. Cargue combustible y busque fugas.

\section{Indicación de cantidad}

\begin{precaucion}
No gire el espárrago central ni la tuerca base del emisor. \textbf{No aplique tensión del sistema a un emisor instalado} para probarlo.
\end{precaucion}

El emisor se instala con su junta, pernos en cruz con torque especial. Para verificarlo, se monta simulando su posición, se coloca el brazo en las cuatro posiciones de la figura y se mide la resistencia con multímetro. Si está fuera de tolerancia, doble el brazo del flotador \emph{hacia arriba} si la resistencia es excesiva y \emph{hacia abajo} si es escasa.

\widefigure[0.85\linewidth]{assets/mm/emisor-combustible.png}{Verificación de los emisores de cantidad: resistencias en tope inferior (0--0,5~$\Omega$), posiciones intermedias y tope superior (90 $\pm$ 2~$\Omega$) para los emisores de vejiga D252-3/-4 y de aluminio A550-1/B122-1. Fuente: MM, figura 12-3.\label{fig:emisor-comb}}

\textbf{Verificación de la indicación}: desabastezca, cierre la válvula y cargue en el tanque principal exactamente la cantidad de la tabla; espere 5 minutos y verifique fugas. Con la batería conectada, el indicador debe marcar entre \textbf{media y una y media ancho de aguja por debajo de 1/4 de tanque}. Si no, verifique el emisor y, si está bien, ajuste o cambie el indicador.

\begin{center}
\begin{tabular}{lc}
\toprule
Configuración & Carga en el principal \\
\midrule
Vejigas principal y auxiliar & 6,2~gal \\
Solo vejiga principal & 5,6~gal \\
Aluminio principal y auxiliar & 6,3~gal \\
Solo aluminio principal & 5,4~gal \\
\bottomrule
\end{tabular}
\end{center}

\subsection{Interruptor de bajo combustible}

Se cambia con el tanque principal fuera, sacándolo por la boca de llenado con dedos mecánicos limpios (sin atrapar el flotador) y guantes sin pelusa; al instalar, junta tórica nueva con A257-6, cables pasados con un alambre de 30~in, sellador B270-1 en la tuerca AN924-5D y torque especial. \textbf{Prueba}: con batería conectada, empuje suavemente el flotador con una varilla de madera limpia: la luz de bajo combustible debe encenderse tras \textbf{aproximadamente 1 segundo}. Si no enciende con el flotador pero sí con el botón de prueba, reemplace el interruptor.

\section{Gascolador}

Limpieza del vaso: batería y válvula cerradas, faldones izquierdos fuera, aeronave conectada a tierra, drenar en un recipiente conectado a tierra. Corte el alambre del collar, retire collar, anillo y vaso; inspeccione y limpie vaso y malla. Al armar: malla y junta en la parte superior; grasa A257-6 en la rosca y el reborde interior del collar; diámetro mayor del anillo contra el collar y menor en la ranura del vaso; apriete a mano hasta asentar \textbf{sin roscas visibles}; frene con alambre de 0,032~in. Haga prueba en tierra y de motor y verifique ausencia de fugas.

\section{Prueba de flujo mínimo de combustible}

\begin{pasos}
\item Batería y válvula cerradas, aeronave a tierra, faldones derechos fuera, tapas puestas.
\item Cargue en el tanque principal exactamente \textbf{3,0~gal} (vejigas) o \textbf{1,7~gal} (aluminio).
\item Desconecte la manguera B283 del carburador, abra batería y válvula y drene a un recipiente conectado a tierra \textbf{hasta que se encienda la luz de bajo combustible}; cierre.
\item Con un recipiente calibrado, abra la válvula y mida el tiempo de llenado:
\end{pasos}

\begin{center}
\begin{tabular}{lc}
\toprule
Cantidad & Tiempo máximo \\
\midrule
1 galón EE.\,UU. & 2~min 25~s \\
4 litros & 2~min 33~s \\
1 galón imperial & 2~min 54~s \\
\bottomrule
\end{tabular}
\end{center}

Si se excede, verifique respiraderos y busque obstrucciones en el filtro de salida, líneas y mangueras. Reconecte la manguera con torque especial, haga prueba de motor y verifique fugas.

\section{Control ambiental}
\fuente{MM, capítulo 11}

Hay ventilaciones de aire fresco en cada puerta y en la nariz. Las de las puertas se abren con la perilla junto a la bisagra (en aeronaves antiguas, empujar la perilla sella y traba; en las recientes hay una perilla giratoria que permite dejarlas parcialmente abiertas). Para máxima ventilación, ábralas del todo en estacionario pero \textbf{solo una pulgada o menos en crucero}. La toma de la nariz se abre tirando de la perilla de aire de cabina de la consola y dirige el aire por el parabrisas para desempañarlo.

\textbf{Ventilación de la nariz}: control A522-1 con contratuerca en la consola, pasador con chaveta nueva en la aleta de la caja; abrazaderas AN742-3 de forma que la funda no resbale; holgura bajo la perilla de \textbf{0,03--0,10~in} con la perilla hundida (cerrada). (Requiere inclinar la consola superior y, si lo hay, retirar el PFD Aspen; puede requerir retirar la batería.)

\textbf{Calefacción de cabina (opcional)}: toma de aire en el lado izquierdo de la voluta del ventilador del motor, camisa de calefacción en el silenciador, válvula en el cortafuego, rejilla delante del asiento del piloto o de los pedales y conductos. Un control tipo empuje-tracción delante del asiento del piloto envía el calor a la cabina o por la borda bajo la cabina. Los calefactores antiguos usaban un soplador eléctrico en el lado izquierdo del compartimiento del motor (disyuntor HEATER, 10~A).

\begin{nota}
En los calefactores antiguos con soplador eléctrico, la camisa puede retirarse en épocas cálidas para prolongar la vida del silenciador: se quitan soplador, camisa y mangueras, y se saca y rotula el disyuntor HEATER. Los calefactores actuales enfrían constantemente el silenciador con aire de la voluta y deben permanecer instalados.
\end{nota}
